% !TeX root = ../main.tex
% !TeX encoding = UTF-8

\chapter{The Fiduciary Standard You Do Not Know You Are Violating}

The letter from the shareholder's counsel runs to four pages. It is addressed to each director individually, by name and title, and references each director's appointment date and the specific board resolutions they voted to approve in the preceding two years. It does not allege fraud. It does not allege gross negligence. It alleges something more difficult to defend against: a failure of informed oversight.

The allegation, stated precisely, is that each named director approved the deployment and continued operation of an AI system in a material decision process without taking the steps necessary to satisfy themselves that the information available to the board at the time of each relevant decision was adequate to discharge their duty of care. The letter does not say the AI system was wrong. It does not say the decisions it contributed to were unfair. It says the directors cannot demonstrate they had the information they needed to govern it, and that, under section 174 of the Companies Act 2006, that failure is actionable.

The directors' first call is to their D\&O insurer. The second is to leading counsel. The third, too late, is to anyone who can explain what section 174 actually requires.

\begin{center}\rule{0.5\linewidth}{0.5pt}\end{center}

The Law You Already Owe, Before Any AI Regulation Applies

This chapter does not provide legal advice. It does something more immediately useful: it identifies the legal obligations that already apply to every company director in England and Wales when they govern an organisation that uses AI in material decisions, before a single AI-specific regulation is considered.

The foundational duties of company directors are not a recent innovation. The duty of care and the duty to act in good faith in the best interests of the company have been features of English law for more than a century. The Companies Act 2006 codified and in some respects tightened those duties. What has changed is not the law. What has changed is the nature of the decisions the law now applies to.

When a board approves the deployment of an AI system in a material decision process, credit, underwriting, hiring, procurement, clinical triage, it has made a decision with consequences it cannot fully anticipate, using a tool it may not fully understand, in a regulatory environment that is actively evolving. The fiduciary duties that apply to that board are the same duties that applied to their predecessors governing paper-based decision processes. They have not been updated for AI. They do not need to be. They are already sufficient to create the liability this chapter describes.

The law governing directors' duties was not written for artificial intelligence. It does not need to have been. It applies to every consequential decision a board governs, regardless of the mechanism that produced it.

\begin{center}\rule{0.5\linewidth}{0.5pt}\end{center}

The Companies Act 2006: Three Duties, One Standard

Three sections of the Companies Act 2006 are directly relevant to any board governing AI-assisted decision processes. Together, they create a standard of conduct that most boards have not yet mapped against their AI governance position.

Section 172 The Duty to Promote the Success of the Company

Section 172 requires a director to act in the way they consider, in good faith, would be most likely to promote the success of the company, including having regard to the long-term consequences of decisions and the need to maintain a reputation for high standards of business conduct. Its relevance to AI governance is frequently underestimated: an AI system producing unfair, biased, or legally vulnerable decisions does not damage the company at the moment of each output. The damage accumulates in the pattern, in enforcement risk, litigation exposure, and reputational harm that materialises when the pattern is examined. Section 172 asks whether a director had sufficient information about the long-term consequences of an AI system's decisions when they approved its deployment and at each point they allowed it to continue operating. Without a Digital Alibi architecture, the honest answer is that they did not, and cannot demonstrate that they did.

Section 174 The Duty of Care, Skill and Diligence

Section 174 requires a director to exercise the care, skill and diligence that would be exercised by a reasonably diligent person with the general knowledge, skill and experience that may reasonably be expected of a person carrying out the functions carried out by the director in relation to the company. This is the duty that the shareholder's counsel invoked in the letter with which this chapter opened.

The standard is a dual standard: objective and subjective simultaneously. The objective component sets a floor, the level of care that any reasonable director carrying out those functions should exhibit. The subjective component raises that floor for directors who hold themselves out as having particular expertise, including expertise in technology, AI, or digital governance. A NED who sits on a technology committee, who has listed AI governance credentials in their biography, or who has led the board's engagement with AI risk, is held to a higher standard of care than the objective floor, because they have implicitly represented a level of knowledge that section 174 requires them to exercise.

For AI governance, section 174's practical requirement is this: a director exercising the requisite care, skill and diligence must be in a position to form an informed view of the material decisions that AI systems are contributing to. An informed view is not a policy-level understanding. It is the capacity to ask and receive answers to specific questions about specific decisions. A director who has approved an AI system in good faith but cannot access information about what that system is actually doing, at the level of individual material decisions, is not exercising section 174 care. They are relying on the assumption that things are working as intended, and reliance on assumption is not diligence.

The Post Office Horizon directors provide the most instructive example of assumption-as-governance. The board's stated position, that the Horizon system was reliable, was not a deliberate misrepresentation. It reflected a genuine institutional belief, held and communicated in good faith. But section 174 does not ask whether a director believed the system was working. It asks whether a reasonably diligent director, with access to the information that was available in the organisation, would have had grounds for that belief. The Horizon Inquiry has progressively established that the information available in the organisation, bug reports, second opinions, legal correspondence, was not the information the board received. The board's belief was genuine. Their access to the information that would have substantiated or challenged it was not. That is the section 174 failure: not dishonesty, but inadequate information access at the moment the function was being exercised.

A director who cannot access information about what their AI system is doing, at the level of individual material decisions, is not exercising the duty of care. They are relying on assumption. Reliance on assumption is not diligence.

Section 175 The Duty to Avoid Conflicts of Interest, A Board Alert

Section 175 surfaces a specific risk that boards should note as AI governance advisory services and AI system vendors increasingly overlap. A director who receives AI governance assurance from a party with a financial interest in a favourable assessment is in a position this duty is designed to address. As the market for AI governance advice matures, the board should satisfy itself that the assurance it receives on AI governance adequacy comes from sources without a financial stake in the conclusion.

\begin{center}\rule{0.5\linewidth}{0.5pt}\end{center}

The SM\&CR: From Corporate Obligation to Personal Jeopardy

Chapter 1 introduced the FCA's Senior Managers and Certification Regime and the specific personal liability mechanism created by section 66A of the Financial Services and Markets Act 2000. Chapter 3 is the place to be precise about what that mechanism requires and, specifically, what ``reasonable steps'' means in an AI governance context.

The SM\&CR operates on a principle of named individual accountability. For every material function within a regulated firm, there must be a named Senior Manager who is accountable for that function and who can, if called upon, demonstrate what steps they took to ensure the function was carried out appropriately. The regime does not operate on team accountability, departmental accountability, or organisational accountability. It operates on named personal accountability, and the word ``named'' is not ceremonial. The Senior Manager must be identified, their scope of responsibility must be documented, and that documentation must be current.

The concept of ``reasonable steps'' under section 66A is not defined in the statute. It has been developed through FCA supervisory practice and a small number of enforcement decisions. The steps that courts and the FCA have treated as reasonable in equivalent contexts include: being present in the governance process rather than merely represented in it; receiving information that is specific enough to allow genuine review rather than general enough to allow easy ratification; having genuine authority to intervene when a concern arises; and creating a contemporaneous record of having done so. What the FCA has never accepted as reasonable steps is the assertion, made after a breach has occurred, that the Senior Manager believed the governance was adequate. Belief, however sincere, is not a step.

In an AI governance context, the reasonable steps a named Senior Manager should take are specific: understanding which AI systems fall within their accountability; ensuring those systems are operated under governance frameworks that capture the Digital Alibi elements at the output moment; being in a position to access information about specific decisions at the level of the output moment rather than aggregate performance metrics; and maintaining a contemporaneous record of their oversight activity. An AI governance programme that provides a Senior Manager with quarterly model performance dashboards, annual validation reports, and a governance policy document does not provide the evidence of reasonable steps that section 66A requires. It provides the appearance of oversight. The SM\&CR requires the substance.

\textbf{WHAT ``REASONABLE STEPS'' REQUIRES FROM A NAMED SENIOR MANAGER IN AI GOVERNANCE}

\begin{itemize}
\tightlist
\item
  Named accountability: your name, your scope, documented and current, not assigned to your role or your team.
\item
  Specific information access: the capacity to interrogate specific AI-assisted decisions at the output moment level, not aggregate performance summaries.
\item
  Genuine intervention authority: a documented mechanism by which you can halt, modify, or override AI system operation, and a record of having used or tested it.
\item
  Contemporaneous record: evidence, created at the time, of each step you took to discharge your oversight responsibilities. Not a retrospective account of what you believe you did.
\end{itemize}

Without all four, a section 66A defence is an assertion. The FCA has repeatedly made clear that assertions are not evidence of reasonable steps.

\begin{center}\rule{0.5\linewidth}{0.5pt}\end{center}

\section{The EU AI Act: The Floor, Not the Ceiling}

The EU AI Act, in force for high-risk AI system operators from 2 August 2026, is the most specific and enforceable AI governance standard yet enacted. Its Articles 9, 12, 13, 14, and 19 create obligations that represent the closest regulatory approach yet to a Digital Alibi mandate. The following table shows precisely where each article reaches, and where it stops.

\begin{longtable}[]{@{}ll@{}}
\toprule\noalign{}
Article & What It Requires \\
\midrule\noalign{}
\endhead
\bottomrule\noalign{}
\endlastfoot
Article 9 & Risk management \\
Article 12 & Record keeping \\
Article 13 & Transparency \\
Article 14 & Human oversight \\
Article 19 & Log retention \\
\end{longtable}

The pattern is identical across all five articles: the Act governs the system, not the output moment. Full compliance builds the conditions under which a Digital Alibi could exist. It does not build the Digital Alibi. The gap is at the layer below, the individual output moment, the specific information picture, the named reviewer, the contemporaneous accountability record. The Act is the floor. The Digital Alibi is what standing on that floor, at the moment each decision was made, actually looks like.

EU AI Act compliance builds the house. The Digital Alibi is the evidence that you lived in it, at the moment each decision was made.

\begin{center}\rule{0.5\linewidth}{0.5pt}\end{center}

\section{The Central Argument: Fiduciary Duty Is Temporal}

Everything in this chapter, the Companies Act duties, the SM\&CR personal liability mechanism, the EU AI Act compliance floor, leads to the argument that is this book's most original contribution to the governance literature. It is an argument that has not been made, to my knowledge, in the published academic or practitioner governance canon. It is this:

Fiduciary duty is not satisfied retrospectively. A director's duty of care is assessed at the moment they made the decision, not at the moment they are questioned about it. The information picture must have existed and been accessible at the decision point. Infrastructure built after the fact does not satisfy the question that arrives before it.

The temporal dimension of fiduciary duty is not a novel legal principle. It is embedded in the Companies Act's language, ``a director must exercise the care, skill and diligence that would be exercised by a reasonably diligent person'', which uses the present tense for a reason. The standard applies at the moment of the exercise of the function, not at the moment of evaluation. A director who makes a decision without adequate information is in breach of section 174 at the moment of that decision, regardless of whether the breach is ever identified, and regardless of what information becomes available subsequently.

Applied to AI governance, the temporal principle produces a specific and uncomfortable conclusion: a board that approved the deployment of an AI system in a material decision process and has not subsequently built a Digital Alibi architecture is not in breach of their duties now. They were in breach at every output moment, for every material decision, at which the Digital Alibi elements were not present. The breach is historical. It is accumulated. And it is not cured by building the architecture today.

This is why the window before regulatory enforcement arrives is not a deadline for becoming compliant. It is a deadline for stopping the accumulation of historical breach. Every output moment that passes without a Digital Alibi record is an output moment at which the board's fiduciary obligations were not met at the information level the law requires. The governance architecture built today does not retroactively cure those moments. It prevents their future recurrence. The accumulation stops when the architecture starts.

\begin{center}\rule{0.5\linewidth}{0.5pt}\end{center}

\section{Three Tiers of Exposure: Where Does Your Board Sit?}

Boards governing AI-assisted decision processes occupy one of three positions relative to the fiduciary standard this chapter has described. The position is determined not by intent or by the quality of AI deployment but by governance architecture, specifically, by the presence or absence of the Digital Alibi elements at the output moment for material decisions.

\begin{longtable}[]{@{}ll@{}}
\toprule\noalign{}
Governance Position & Description \\
\midrule\noalign{}
\endhead
\bottomrule\noalign{}
\endlastfoot
Tier 1 & No formal AI governance \\
Tier 2 & Voluntary commitments only \\
Tier 3 & Formal governance active, Digital Alibi gap remains \\
\end{longtable}

The majority of regulated organisations deploying AI in material decisions today sit at Tier 2 or Tier 3. Tier 1 organisations are increasingly uncommon in regulated sectors, though they are prevalent in mid-market and PE-backed companies without formal compliance functions. Tier 3 is the most dangerous position, not because the governance is inadequate by conventional standards, but because it creates the expectation of adequacy without the substance that a challenge will require.

A Tier 3 board that faces a sophisticated regulatory inquiry or a shareholder challenge will produce its governance documentation with confidence. The regulator will examine it, note that it is thorough, and ask the output-moment question. The governance documentation will not answer the output-moment question. The board will then attempt to reconstruct the output-moment evidence from logs, emails, and model documentation, and the reconstruction will be examined for the gap between what it claims and what the contemporaneous record would have shown. The Tier 3 board's confidence was not misplaced. It was simply applied to the wrong question.

\begin{center}\rule{0.5\linewidth}{0.5pt}\end{center}

\section{A Worked Example: Insurance Underwriting at Lloyd's}

A Lloyd's managing agent deploys an AI system to assist underwriters in pricing complex specialty risks. The system is sophisticated, procured from a specialist vendor with a strong market reputation, and has been deployed through a formal technology governance process that includes model validation, risk committee sign-off, and board-level approval. The managing agent has an SM\&CR accountability map that names the Chief Underwriting Officer as the Senior Manager responsible for underwriting decisions. The board receives quarterly reports on underwriting performance, including model performance statistics.

Eighteen months after deployment, Lloyd's Corporation issues a supervisory notice to the managing agent following a pattern of claims that suggests systemic underpricing in a specific specialty class. The notice asks the managing agent to demonstrate its governance position for the AI-assisted underwriting decisions made in the relevant class during the preceding twelve months.

The questions Lloyd's asks are precise: for each of the twenty most material underwriting decisions in the class during the period, produce the information picture that existed at the moment each risk was priced. This means the version of the AI model deployed at each pricing moment; the exposure data inputs presented to the model for each specific risk; the data the model was not given, third-party catastrophe model outputs, for instance, that were available but not integrated; the underwriter's review record for each pricing output; and the accountability chain from model to board current at each decision moment.

The managing agent cannot produce these items for any of the twenty decisions, because its governance architecture captured performance data in aggregate, not evidence at the output moment. The quarterly board reports contained model performance metrics across the specialty class as a whole. They did not contain the information picture for each individual pricing decision. The SM\&CR accountability map named the Chief Underwriting Officer, but the map was created at the point of the system's deployment. Three of the twenty decisions were made in a period when the previous Chief Underwriting Officer had resigned and their successor had not yet been formally appointed to the SM\&CR function. For those three decisions, there was no named accountable Senior Manager at the output moment.

The Lloyd's supervisory notice becomes a formal investigation. The managing agent's defence, that it had a sound governance framework, properly approved by the board, with appropriate model validation and SM\&CR documentation, is accurate. It is also irrelevant to the questions the investigation is asking. The governance framework describes what the managing agent intended to do. The investigation is asking what it did, at the moment of each specific pricing decision. These are different questions. The framework cannot answer the second.

The governance framework describes what was intended. The investigation asks what was done, at the moment of each specific decision. A framework cannot answer for moments it was not designed to capture.

\begin{center}\rule{0.5\linewidth}{0.5pt}\end{center}

\section{The Question Boards Should Be Asking}

The fiduciary obligation described in this chapter does not require boards to understand how AI systems work. It does not require technical expertise in machine learning, statistical modelling, or software architecture. It requires boards to be able to answer one question, for any material AI-assisted decision process they have approved: at the output moment, for each governed decision, are the four Digital Alibi elements present in a contemporaneous, retrievable form?

That question should be on every board's AI governance agenda. Not as a standing item about AI strategy, AI risk appetite, or AI ethics. As a specific, operational question about evidence. The board does not need to examine individual decision records. It needs to know that the architecture exists to produce those records, that the output-moment capture is operational, that the accountability chains are current, that the human oversight records are genuine rather than performative, and that the retrieval mechanism has been tested.

\textbf{THE FOUR QUESTIONS A BOARD SHOULD RECEIVE ANSWERS TO AT EVERY AI GOVERNANCE REVIEW}

\begin{enumerate}
\def\labelenumi{\arabic{enumi}.}
\tightlist
\item
  For each material AI system, is output-moment capture operational, producing a contemporaneous record of the information picture for each decision as it occurs?
\item
  Are the accountability chains current, with named individuals documented at the decision-moment level, not assigned to roles or updated only at annual review?
\item
  Are the human oversight records genuine, with evidence that named reviewers had the information, time, and authority to override, rather than a sign-off process that is structurally incapable of producing overrides?
\item
  Has the retrieval mechanism been tested, can we produce the complete Digital Alibi for a specific decision, on demand, within the timeframe a regulator or court would require?
\end{enumerate}

A board that cannot answer yes to all four questions has a Digital Alibi gap. The gap may be small, a retrieval mechanism that has never been tested, an accountability chain that is current for eleven months and outdated for one, or it may be structural, representing a governance architecture that has never been designed to capture output-moment evidence. Either way, the gap is a fiduciary exposure. The law does not grade on a curve between a small gap and a structural one. The duty of care either was or was not met at each output moment. The digital alibi either exists or it does not.

\begin{center}\rule{0.5\linewidth}{0.5pt}\end{center}

\section{Why Existing Frameworks Do Not Close This Gap}

Chapter 3 has established the legal architecture of the fiduciary exposure the Digital Alibi gap creates. It has identified the Companies Act duties, the SM\&CR personal liability mechanism, and the EU AI Act as the three primary instruments through which the exposure is structured, and demonstrated that none of them, individually or together, close the gap at the output-moment level.

Chapter 4 examines why this is so, not as a criticism of the frameworks, but as a structural analysis. The AI governance canon that has been built over the past decade is a serious body of work produced by serious people. It addresses process, policy, and system-level governance with genuine rigour. The Digital Alibi gap is not the result of the frameworks being inadequate at what they were designed to do. It is the result of the frameworks not having been designed to address the evidentiary question that a regulator's examination, a shareholder's challenge, or a court's inquiry will actually ask.

Understanding why the frameworks miss this question, and why they were not designed to address it, is essential to understanding what filling the gap requires. Chapter 4 conducts that analysis, framework by framework, and closes Part One of this book with the complete account of the gap: what it is, how it was created, and why it has persisted unaddressed until now.
